\section{实体关系抽取}
\label{sec:nlp-extraction-relations}

找全组织和活动名称以后，还要判断谁举办了什么。\term{关系抽取}（relation extraction, RE）
把文本中的实体与它们之间的语义关系绑定起来，输出有序的关系三元组
\[
  (s,r,o),\qquad r\in R\cup\{\mathrm{NA}\},
\]
其中$s$为\emph{主体}（subject），$o$为\emph{客体}（object），$R$是本体预先给定的关系类型集合，
$\mathrm{NA}$表示两者之间没有$R$中的任何关系。三元组是有序的：
$(\text{海棠大学},\text{主办},\text{读书会})$表示前者主办后者，交换两端就成了另一个命题。
本节用到的组织、人物和句子都是补充的教学材料，不由章首通知推断。
二元关系有两个参与者；“海棠大学于4月18日在海棠楼201主办读书会”这样的多元事实
还带有时间和地点。若把它拆成几条二元三元组存储，需要给它们一个共同的实例键，
让时间、地点都挂到同一次“主办”上；否则另一句“晨光书院于4月19日在松柏楼举办研讨会”
的日期和地点可能被拼进这条记录。多元事实的完整处理留给第~\ref{sec:nlp-extraction-events}节的事件抽取，
事件实例之间的先后与因果见第~\ref{sec:nlp-extraction-event-relations}节。

句内或文档级说明允许读取的上下文，封闭模式或开放关系说明标签是否预先枚举，
二元或多元说明参与结构。它们是可以组合的任务条件，不是互斥的模型类别。

\subsection{关系抽取的挑战与三元组重叠类型}
\label{sec:nlp-extraction-relation-challenges}

关系抽取的困难可以归为五项。其一是候选规模：一个文本中有$M$个实体时，
不含自关系的有向实体对有$M(M-1)$个，4个实体就有12对，其中绝大多数是NA，
类别极不平衡。其二是三元组之间共享实体，一个实体可能同时参与多条关系，
逐词元只给一个标签的输出表示装不下这种结构。其三是监督稀缺：逐句标注关系代价高，
借助知识库自动对齐得到的标签又带有大量噪声。其四是证据跨句：一个实体在文档中有多个提及，
关系可能要结合几句话才能判断。其五是类型不封闭：预先定义的$R$总会遗漏新出现的关系表达。
后续各小节分别回应这些困难：流水线方法和PURE处理候选规模与实体信息的注入，
远程监督与多示例学习处理监督稀缺，联合抽取处理实体共享，文档级方法处理跨句证据，
开放关系抽取处理不封闭的类型。

第二项困难需要先说清“共享”有哪几种。设句子的金标准三元组集合为$T$，
记$\operatorname{ent}(t)=\{s,o\}$为三元组$t=(s,r,o)$涉及的两个实体。
\citeauthor{zeng2018copyre}按实体共享方式把句子分为三类\cite{zeng2018copyre}：
\begin{itemize}
  \item \emph{无重叠}（Normal）：任意两条三元组不共享实体，即对所有$t\neq t'$有
    $\operatorname{ent}(t)\cap\operatorname{ent}(t')=\varnothing$；
  \item \term{实体对重叠}（entity pair overlap, EPO）：存在两条三元组共享同一对实体，
    即$\operatorname{ent}(t)=\operatorname{ent}(t')$而关系不同；
  \item \term{单实体重叠}（single entity overlap, SEO）：存在两条三元组恰好共享一个实体，
    即$|\operatorname{ent}(t)\cap\operatorname{ent}(t')|=1$。
\end{itemize}
同一句话可以同时属于EPO和SEO。图~\ref{fig:nlp-extraction-overlap}给出三类句子的实例。
“林舟出生并工作于海棠市”有$(\text{林舟},\text{出生地},\text{海棠市})$与
$(\text{林舟},\text{工作地},\text{海棠市})$，两条三元组的实体对完全相同，属于EPO。
“海棠大学和晨光书院共同主办读书会”有两条“主办”三元组，主体不同、客体都是读书会，
属于SEO；“林舟创办的晨光书院位于海棠市”中，晨光书院在一条三元组里是客体、
在另一条里是主体，也属于SEO。两个主体共享客体只说明它们各自主办该活动，
并不蕴含两个主体之间存在“合作”关系；后者若需要，应作为另一条有原文依据的三元组。

\begin{figure*}[htbp]
  \centering
  % 三元组重叠的三种类型：Normal、EPO、SEO。节点为实体，有向边为关系，边上标注关系类型。
\begin{tikzpicture}[
  x=1mm,y=1mm,figure picture,
  ent/.style={figure node,minimum height=7mm,minimum width=13mm,
    draw=FigureInput,fill=FigureInputFill,font=\figurefont\scriptsize},
  shared/.style={ent,draw=FigureOutput,fill=FigureOutputFill,line width=1pt},
  rel/.style={figure flow},
  rel label/.style={figure edge label,font=\figurefont\scriptsize},
  sent/.style={font=\figurefont\scriptsize,text=FigureMuted,align=center}
]
  % 三个分区
  \figurePanel{0}{0}{44}{40}{Normal}{FigurePanel}
  \figurePanel{50}{0}{44}{40}{EPO}{FigurePanel}
  \figurePanel{100}{0}{48}{40}{SEO}{FigurePanel}

  % Normal：两个三元组不共享实体（此处只画一条）
  \node[ent] (n1) at (9,19) {海棠大学};
  \node[ent] (n2) at (35,19) {读书会};
  \draw[rel] (n1) -- node[rel label,above=1mm] {主办} (n2);
  \node[sent] at (22,6) {海棠大学主办读书会};

  % EPO：同一有序实体对上有两种关系
  \node[shared] (e1) at (59,19) {林舟};
  \node[shared] (e2) at (85,19) {海棠市};
  \draw[rel] (e1.north east) to[bend left=28] node[rel label] {出生地} (e2.north west);
  \draw[rel] (e1.south east) to[bend right=28] node[rel label] {工作地} (e2.south west);
  \node[sent] at (72,4) {林舟出生并工作于海棠市};

  % SEO：两个三元组只共享一个实体
  \node[ent] (s1) at (110,26) {海棠大学};
  \node[ent] (s2) at (110,12) {晨光书院};
  \node[shared] (s3) at (139,19) {读书会};
  \draw[rel] (s1.east) -- node[rel label,above=.5mm,pos=.45] {主办} (s3.north west);
  \draw[rel] (s2.east) -- node[rel label,below=.5mm,pos=.45] {主办} (s3.south west);
  \node[sent] at (124,4) {海棠大学和晨光书院共同主办读书会};
\end{tikzpicture}

  \caption{三元组重叠的三种类型。方框为实体，有向边从主体指向客体，边上为关系类型，
  加粗方框为被多条三元组共享的实体，方框下方为对应的教学句。无重叠的句子中各三元组互不共享实体；
  EPO中同一对实体之间有两种关系；SEO中两条三元组只共享一个实体（此处为共同的客体）。
  三类划分依据Zeng等（2018）的定义，例句为教学构造。}
  \label{fig:nlp-extraction-overlap}
\end{figure*}

重叠决定了输出表示的最低要求。若每个词元只能带一个标签，EPO中的实体对要同时表达两种关系，
SEO中的共享实体要同时扮演两个角色，都装不下。若改为对每个有序实体对分别判断，
SEO自然可以表示；EPO则还要求同一实体对允许多个关系，即多标签而不是单标签输出。
表~\ref{tab:nlp-extraction-relation-methods}按这一标准对照本节的方法，
各行的具体机制在后续小节展开。

\begin{table*}[htbp]
  \centering
  \small
  \begin{tabularx}{\linewidth}{@{}>{\raggedright\arraybackslash}p{24mm}
    >{\raggedright\arraybackslash}X
    >{\centering\arraybackslash}p{16mm}>{\centering\arraybackslash}p{8mm}
    >{\raggedright\arraybackslash}X@{}}
    \toprule
    方法 & 输出表示与核心机制 & EPO & SEO & 主要代价 \\
    \midrule
    实体对分类（流水线、PURE） & 先识别实体，再对每个有序实体对分类 & 多标签时能 & 能 & 实体错误向下传播；每对单独编码 \\
    PCNN、选择性注意力 & 同一实体对的句子组成包，包级分类 & 多标签时能 & 能 & 依赖远程监督假设；只处理给定实体对 \\
    NovelTagging & 每词一个“位置—关系—角色”组合标签 & 否 & 否 & 无法表示重叠 \\
    CopyRE & 每三步生成关系并复制两端实体 & 能 & 能 & 只复制实体末词；三元组数目与次序固定 \\
    CasRel & 主体首尾指针，再按关系条件标客体首尾 & 能 & 能 & 主体漏检丢失其全部关系；训练与推断条件不一致 \\
    表格填充 & $n\times n$词对表，每格一个标签 & 否 & 能 & 一格只能存一种关系；填表顺序影响结果 \\
    TPLinker & 每种关系一组首首、尾尾词元对连接 & 能 & 能 & 输出规模$O(|R|n^2)$ \\
    GAIN、ATLOP、EIDER & 文档级实体对多标签分类 & 能 & 能 & 实体与提及须先给定；长文档编码 \\
    ReVerb & 词性模式匹配关系短语，再找左右论元 & 能 & 能 & 规则面向英语；关系不归一 \\
    \bottomrule
  \end{tabularx}
  \caption{关系抽取方法谱系。EPO与SEO两列表示输出表示能否容纳该重叠类型，
  不表示实际抽取的准确率。“多标签时能”指同一实体对允许输出多个关系时才能表示EPO。}
  \label{tab:nlp-extraction-relation-methods}
\end{table*}

\subsection{流水线关系抽取}
\label{sec:nlp-extraction-relation-pipeline}

最直接的做法分两步。第一步用第~\ref{sec:nlp-extraction-entities}节的实体识别得到实体集合
$\widehat{E}$；第二步对每个有序实体对$(e_i,e_j)$，$i\neq j$，计算
\[
  \hat r_{i,j}=\argmax_{r\in R\cup\{\mathrm{NA}\}}p(r\mid x,e_i,e_j).
\]
这种\emph{流水线}（pipeline）结构的好处是两个模型可以分别设计、分别训练，
实体识别还能借用现成的标注数据和模型。局限也来自分步。

第一个局限是\term{误差传播}（error propagation）：一条三元组正确，要求两个实体的边界与类型都对，
关系也对。设实体识别对每个金标准实体的召回率为$\rho_e$，两端识别近似独立，
关系分类器在实体正确的前提下判对的概率为$\rho_r$，则三元组召回率约为$\rho_e^{2}\rho_r$。
取教学数值$\rho_e=\rho_r=0.9$，三元组召回率只剩$0.9^{3}\approx0.729$；
第一步漏掉的实体，第二步无论多准都无法补回。错误的实体还会制造虚假候选，
降低精确率。第二个局限是训练与推断的条件不一致：关系模型通常在金标准实体上训练，
推断时面对的却是带错误的预测实体。第三个局限是代价：每个实体对若要单独编码，
一句话有12个候选对就要做12次前向计算。联合抽取（第~\ref{sec:nlp-extraction-joint}小节）
针对第一、二项，用一个模型同时输出实体和关系；PURE则说明，只要把实体信息正确地注入关系模型，
分开的两个模型仍有很强的竞争力，并给出处理第三项的近似计算。

\subsubsection{句内实体对分类与依存路径特征}
\label{sec:nlp-extraction-sentence-re}

给定句子和一个有向提及对$(m_s,m_o)$，关系分类要回答：这句话是否表达了$m_s$到$m_o$的
某种关系。早期方法手工构造特征向量$\bm{\phi}$，包括两个实体的词与类型、两者的相对距离、
两者之间的词，以及两实体在依存树上的\term{最短依存路径}（shortest dependency path）。
\citeauthor{bunescu2005shortest}的出发点是：依存树中连接两个实体的最短路径集中了判断关系所需的
大部分信息，路径以外的修饰成分多半与关系无关\scite{bunescu2005shortest}
依存树的构造见第~\ref{sec:nlp-analysis-dependency}小节。
补充句“海棠大学主办读书会”的路径为
“海棠大学$\xleftarrow{\text{nsubj}}$主办$\xrightarrow{\text{obj}}$读书会”，
弧的方向和标签告诉模型谁是施事、谁是受事，这正是区分$(s,r,o)$与$(o,r,s)$的依据。
分数由线性层给出，$\bm{z}=\bm{W}\bm{\phi}+\bm{b}\in\mathbb{R}^{|R|+1}$，
额外一维对应NA，再经softmax得到各关系的概率。

预训练编码器出现后，特征改由模型自己学。常用做法是在输入中插入\term{实体标记}（entity marker）：
在主体前后加\code{[E1]}、\code{[/E1]}，在客体前后加\code{[E2]}、\code{[/E2]}，
编码后取两个起始标记位置的向量拼接，
$\bm{\phi}=[\bm{h}_{\texttt{[E1]}};\bm{h}_{\texttt{[E2]}}]\in\mathbb{R}^{2d}$，
再以$\bm{W}\in\mathbb{R}^{(|R|+1)\times2d}$分类\scite{baldinisoares2019mtb}
标记把“哪两个跨度是当前候选”写进了输入，编码器的每一层都能据此调整上下文表示；
同一句话换一个实体对，就是一个不同的输入。两种标记区分了主客体，
交换实体顺序会得到另一个输入，模型不默认关系对称。
单标签设定用softmax交叉熵训练；若同一对允许多个关系（EPO），则改为逐关系的二分类，
或第~\ref{sec:nlp-extraction-atlop}小节的自适应阈值目标。

举一个教学数值。设三个类别“主办、支持、NA”的得分为$(2,0,-1)$，则
\[
  p(\text{主办})=\frac{\mathrm{e}^{2}}{\mathrm{e}^{2}+\mathrm{e}^{0}+\mathrm{e}^{-1}}
  \approx\frac{7.389}{8.757}\approx0.8438,
\]
正确标签的负对数似然约为0.1698，解码输出有向的“主办”关系。
依存路径特征的弱点可以用一个反例看出：若句子变成“海棠大学并不主办读书会”，
“并不”修饰“主办”，却不在两实体之间的最短路径上，只看路径的特征与原句完全相同，
分类器会给出同样的“主办”。路径之外的否定、情态和条件修饰必须另作特征或进入编码器表示；
带实体标记的编码器读入整句，可以利用这些成分，但是否学会利用仍取决于训练数据中有无此类样本。

\subsubsection{PURE实体标记与近似加速}
\label{sec:nlp-extraction-pure}

联合抽取的一个常见论点是：共享编码器能让实体和关系互相提供信息，从而克服流水线的缺陷。
\citeauthor{zhong2021pure}提出的PURE反其道而行，用两个\emph{不共享}参数的编码器分别做实体识别和关系分类，
结论是关系模型的关键在于尽早获得实体的边界与类型\scite{zhong2021pure}
它的实体模型是跨度分类器：枚举宽度不超过$W$的跨度$(i,j)$，以
$\bm{g}_{i,j}=[\bm{h}_i;\bm{h}_j;\bm{\phi}(j-i)]$表示，即式~\eqref{eq:nlp-analysis-span-representation}
去掉内部注意力汇聚的版本，再经前馈网络在“实体类型$\cup\{\varepsilon\}$”上分类，
$\varepsilon$表示不是实体。为了利用跨句信息，输入在句子左右各扩展若干上下文词元。

关系模型把实体模型的输出写进输入。对一个候选对，在主体两端插入带类型的标记
\code{<S:ORG>}、\code{</S:ORG>}，在客体两端插入\code{<O:EVT>}、\code{</O:EVT>}，
这里ORG、EVT是实体模型预测的类型。例如
\[
  \texttt{<S:ORG>}\ \text{海棠大学}\ \texttt{</S:ORG>}\ \text{主办}\
  \texttt{<O:EVT>}\ \text{读书会}\ \texttt{</O:EVT>}.
\]
编码后取两个起始标记的输出拼接，
$\bm{u}_{s,o}=[\hat{\bm{x}}_{\texttt{<S:ORG>}};\hat{\bm{x}}_{\texttt{<O:EVT>}}]\in\mathbb{R}^{2d}$，
经线性层在$R\cup\{\varepsilon\}$上做softmax，训练时实体取金标准。
与上一小节不带类型的标记相比，类型标记让关系模型直接得知“这是组织对活动”，
可以排除大量类型不相容的关系。若把实体模型的跨度向量拼接到关系分类器上，
实体信息只在最后一层才进入，编码器在计算上下文表示时并不知道当前是哪一对；
PURE的实验对照显示，这种后期融合不如在输入层插入标记有效。

代价是每个实体对都要重新编码一次。设句子有$n=30$个词元、$M=4$个实体，
共12个有序对，精确做法要做12次前向计算，每次长度约34。
PURE的近似做法把所有实体对的标记移到句末，标记与它所标跨度的首、尾词元共享位置编号，
并用注意力掩码规定：原文词元只看原文词元，标记看原文词元和同一实体对的其他标记。
这样原文词元的表示与实体对无关，只需计算一次；每对的四个标记各自读取原文并彼此交互。
一句话只需一次前向，长度为$30+4\times12=78$。
近似的代价是原文表示不再随实体对变化，标记信息不能反过来影响上下文，
论文报告推断加速8到16倍，F1仅小幅下降。

PURE说明，误差传播并不必然使流水线劣于联合模型：实体模型足够准、实体信息又在关系模型
输入层注入时，分开训练的两个模型可以与当时的联合模型相当或更好。但流水线的结构性问题仍在，
实体模型漏掉的实体不会再有关系预测，关系模型也仍在金标准实体上训练，
推断时面对的预测实体类型若有错，带类型的标记反而会把错误类型直接告诉关系模型。

\subsection{远程监督与多示例学习}
\label{sec:nlp-extraction-distant}

逐句标注关系的代价很高。\term{远程监督}（distant supervision）利用已有知识库自动产生训练标签：
若知识库中有三元组$(e_1,r,e_2)$，就把所有同时提及$e_1$与$e_2$的句子都标为关系$r$
\scite{knowledge-acquisition-mintz2009}
这一假设过强。设知识库记有$(\text{海棠大学},\text{主办},\text{读书会})$，句子
“海棠大学主办的读书会改至周五”确实表达了主办，而“海棠大学图书馆张贴了读书会海报”只是共同提及，
也会被标为“主办”。反方向上，知识库不完整，库中没有的关系会让表达了该关系的句子被标为NA。
常用的NYT数据集就是把Freebase中的关系与《纽约时报》语料对齐得到的\cite{knowledge-acquisition-riedel2010}，两类噪声都普遍存在。

\citeauthor{knowledge-acquisition-riedel2010}把假设放宽为“至少一次表达”：
同时提及两实体的句子中，至少有一句表达了库中的关系\scite{knowledge-acquisition-riedel2010}
这正是\term{多示例学习}（multiple-instance learning）的设定：训练数据由若干\emph{包}组成，
每个包是一组示例，标签只给到包，不给到其中的示例；正包中至少有一个正示例，
负包中全部是负示例。监督约束施加在整组对象上，示例级标签始终未知。
这一设定的一般讨论可参考\BookChapterRef{chap:machine-learning-paradigms}{卷三的范式简介章}。
在关系抽取中，同一实体对$(e_1,e_2)$的全部句子组成包$B=\{\bm{x}_1,\ldots,\bm{x}_m\}$，
包标签取知识库中的关系$y$。学习目标是包级的：
\[
  \max_{\bm{\theta}}\sum_{k}\log p_{\bm{\theta}}(y_k\mid B_k),\qquad
  p_{\bm{\theta}}(y\mid B)=f\bigl(\bm{x}_1,\ldots,\bm{x}_m\bigr),
\]
不同方法的区别在于聚合函数$f$怎样从句子得到包：PCNN只选一句，选择性注意力对全部句子加权。
两种方法都以给定实体对为前提，属于流水线的第二步；评价常用留出的知识库三元组作答案，
库中没有的正确预测会被算作错误，因此得到的精确率是保守估计。

\subsubsection{PCNN分段池化}
\label{sec:nlp-extraction-pcnn}

先看句子怎样编码。卷积网络用于关系分类时，常对整句的卷积特征做一次最大池化，
得到与句长无关的定长向量；但一次池化只保留每个滤波器在全句的最强响应，
丢掉了这个响应出现在两个实体之前、之间还是之后。对关系而言，这一位置信息很要紧：
“海棠大学主办读书会”中，动词出现在两个实体之间。
\citeauthor{zeng2015pcnn}的\emph{分段卷积网络}（piecewise convolutional neural network, PCNN）
按两个实体的位置把特征切成三段，分别池化\scite{zeng2015pcnn}

设句子有$n$个词元，主体与客体分别位于位置$p_1<p_2$。第$i$个位置的输入是词嵌入与
两个相对位置嵌入的拼接，
$\bm{x}_i=[\bm{w}_i;\bm{d}^{(1)}_{i-p_1};\bm{d}^{(2)}_{i-p_2}]\in\mathbb{R}^{d_w+2d_p}$，
两个位置嵌入分别编码该词到主体、客体的距离。宽度为$w$的第$k$个滤波器
$\bm{f}_k\in\mathbb{R}^{w(d_w+2d_p)}$在每个位置给出响应
$c_{k,i}=\bm{f}_k^{\top}\bm{x}_{i-w+1:i}+b_k$，
$\bm{x}_{i-w+1:i}$是相邻$w$个输入向量的拼接，句首句尾按需补零。
\term{分段最大池化}（piecewise max pooling）把位置分成三段
$I_1=\{i\leq p_1\}$、$I_2=\{p_1<i\leq p_2\}$、$I_3=\{i>p_2\}$，分别取最大：
\begin{equation}
  q_{k,\ell}=\max_{i\in I_\ell}c_{k,i},\quad \ell=1,2,3,\qquad
  \bm{g}=\tanh\bigl([q_{1,1},q_{1,2},q_{1,3},\ldots,q_{K,3}]\bigr)\in\mathbb{R}^{3K}.
  \label{eq:nlp-extraction-pcnn}
\end{equation}
$K$个滤波器得到$3K$维句子表示，经dropout和线性层在$R\cup\{\mathrm{NA}\}$上做softmax。
原文的设置为$d_w=50$、$d_p=5$、$w=3$、$K=230$，句子表示为690维。
举一个教学数值：$n=7$、$p_1=2$、$p_2=5$，某滤波器的响应为
$(0.2,0.9,0.1,-0.3,0.4,0.6,0.8)$。整句最大池化只得到0.9；
分段池化得到$(0.9,0.4,0.8)$，记录了实体之间也有中等响应、客体之后响应较强，
后续线性层可以据此区分不同的关系模式。

包级训练采用“至少一次表达”的最直接实现：对第$k$个包，先用当前模型找出最可能表达包标签的句子，
\[
  j^{*}=\argmax_{j}p(y_k\mid\bm{x}_{k,j}),\qquad
  J(\bm{\theta})=\sum_{k}\log p(y_k\mid\bm{x}_{k,j^{*}}),
\]
只用这一句计算梯度。例如一个“主办”包有三句，模型给出的$p(\text{主办})$为$(0.2,0.7,0.1)$，
则选第二句，该包的损失为$-\log0.7\approx0.3567$，另外两句不贡献梯度。
推断时，只要包内有一句被判为某关系，就给该实体对输出这一关系。
第~\ref{sec:nlp-extraction-dmcnn}小节事件检测中的DMCNN用的是同一类动态多池化，
切分点换成触发词和候选论元。PCNN的分段池化代价极低，却给卷积编码补上了实体相对位置的结构；
包级只选一句则浪费了其余句子中的有效信息，且训练初期模型选中的句子本身可能就是噪声句。

\subsubsection{句子级选择性注意力}
\label{sec:nlp-extraction-selective-attention}

只选一句的做法有两个问题：一个包中常有多句表达同一关系，其余各句的信息被丢掉；
选择是硬的，一旦选错，梯度全部落到噪声句上。
\citeauthor{lin2016selective}改用\term{选择性注意力}（selective attention），
为包内每句计算一个与关系相关的权重，再加权汇聚成包表示\scite{lin2016selective}
设包内句子表示为$\bm{x}_1,\ldots,\bm{x}_m\in\mathbb{R}^{d_c}$（例如PCNN的$3K$维输出），
关系$r$有查询向量$\bm{r}\in\mathbb{R}^{d_c}$，$\bm{A}$为可学习的对角矩阵，
\begin{equation}
  \begin{aligned}
    e_i&=\bm{x}_i^{\top}\bm{A}\bm{r},\qquad
    \alpha_i=\frac{\exp e_i}{\sum_{j=1}^{m}\exp e_j},\\
    \bm{s}&=\sum_{i=1}^{m}\alpha_i\bm{x}_i,\qquad
    p(r\mid B)=\operatorname{softmax}(\bm{M}\bm{s}+\bm{d})_r.
  \end{aligned}
  \label{eq:nlp-extraction-selective-attention}
\end{equation}
$e_i$衡量第$i$句与关系$r$的匹配程度，$\bm{M}\in\mathbb{R}^{(|R|+1)\times d_c}$是关系分类矩阵。
训练时用包标签对应关系的查询向量，目标是$\sum_k\log p(y_k\mid B_k)$；
推断时对每个候选关系分别用自己的查询计算包表示和分数，因为不同关系应关注不同的句子。

设一个三句包对“主办”的匹配分数为$(1,3,0)$，则
\[
  \bm{\alpha}=\frac{(\mathrm{e}^{1},\mathrm{e}^{3},\mathrm{e}^{0})}
  {\mathrm{e}^{1}+\mathrm{e}^{3}+\mathrm{e}^{0}}\approx(0.114,0.844,0.042),
\]
包表示主要来自第二句，但第一句仍贡献约11\%。这一权重可以看成在两个极端之间：
PCNN的选一句相当于把最大者的权重设为1，其余为0；对全部句子取平均则让噪声句与有效句等权。
注意力是硬选择的平滑版本，梯度能同时流向多个句子，噪声句的权重则随训练降低。
它仍建立在“至少一次表达”的假设上：若包内没有一句真正表达该关系，
注意力也只能在噪声句之间分配权重；知识库缺失造成的假NA包同样无法纠正。
远程监督数据中只含一句的包很多，此时注意力退化为该句本身，聚合带来的好处有限。

\subsection{实体与关系的联合抽取}
\label{sec:nlp-extraction-joint}

\term{联合抽取}（joint extraction）用一个模型同时输出实体与关系，
目的是避免流水线中实体错误无法纠正、实体与关系模型训练条件不一致的问题。
它的难点在于输出表示：既要装下实体边界，又要装下有向关系，还要允许重叠。
下面四种方法依次修补前一方法的缺陷。NovelTagging把三元组编码进一层序列标签，
却装不下重叠；CopyRE改为逐个生成三元组，同一实体可以复制多次，
但只复制实体末词，三元组的数目和次序也被固定；CasRel把关系看作从主体到客体集合的映射，
以首尾指针完整恢复多词实体并处理重叠，但推断依赖先预测出的主体；
表格填充与TPLinker把全部实体与关系同时标在词元对上，一步完成解码。
四种方法沿用同一个SEO教学句，词元依次为
“海棠大学／和／晨光书院／共同／主办／读书会”，金标准为$(1,\text{主办},6)$和$(3,\text{主办},6)$，
编号表示单词元提及。

\subsubsection{NovelTagging统一标注}
\label{sec:nlp-extraction-noveltagging}

把关系抽取改写成序列标注，就可以直接沿用第~\ref{sec:nlp-analysis-paradigm-tagging}小节的标注模型。
\citeauthor{zheng2017noveltagging}为此设计了一套组合标签\scite{zheng2017noveltagging}
每个标签由三部分组成：词在实体中的位置（B、I、E、S），该实体参与的关系类型，
以及它在关系中的角色（1为主体，2为客体）；不属于任何三元组的词标O。标签总数为$2\times4\times|R|+1$。
“海棠大学主办读书会”标为“S-主办-1，O，S-主办-2”。
解码时先按位置标签恢复实体，再把关系类型相同、角色互补的实体配成三元组；
同类实体有多个时，按就近原则配对。
模型是双向LSTM编码器加一个逐位置的LSTM解码器，训练目标在交叉熵上加大非O标签的权重：
\[
  L=-\sum_{t=1}^{n}\Bigl[\mathbb{I}(y_t=\mathrm{O})\log p(y_t\mid x)
    +\alpha\,\mathbb{I}(y_t\neq\mathrm{O})\log p(y_t\mid x)\Bigr],\qquad\alpha>1,
\]
因为O标签占绝大多数，不加权时模型倾向于把关系词也标成O。

统一标注把实体与关系放进同一个输出，一次解码完成，免去了流水线的两步。
它的根本限制在表示上：每个词只有一个标签，一个实体在整句中只能有一种关系、一个角色。
EPO句“林舟出生并工作于海棠市”中，“林舟”需要同时带“出生地-1”和“工作地-1”；
链式SEO句“林舟创办的晨光书院位于海棠市”中，“晨光书院”需要同时带“创办-2”和“位于-1”，
这些都无法用一个标签表示。共享客体的SEO教学句虽然可以标成
“S-主办-1，O，S-主办-1，O，O，S-主办-2”，就近配对却只会把读书会与较近的晨光书院配成一条，
海棠大学的三元组随之丢失。重叠在真实语料中很常见，这一限制促使后续工作改换输出表示。

\subsubsection{CopyRE复制式三元组生成}
\label{sec:nlp-extraction-copyre}

既然一个词元装不下多个角色，就不再给词元贴标签，而是把三元组一条一条地“写”出来：
同一个实体在不同三元组中出现几次，就复制几次。CopyRE用编码器—解码器实现这一想法
\scite{zeng2018copyre}
编码器为双向LSTM，输出$\bm{o}^{E}_1,\ldots,\bm{o}^{E}_n\in\mathbb{R}^{d}$；
解码器每三步生成一条三元组，第$t$步状态为$\bm{o}_t$。三类步骤输出不同的对象：
\begin{equation}
  \begin{aligned}
    &\text{关系步}\ (t\equiv1):&
      p(r)&=\operatorname{softmax}\bigl([\operatorname{selu}(\bm{W}_r\bm{o}_t+\bm{b}_r);\ a_t]\bigr)_r,\\
    &\text{实体步}\ (t\equiv2,0):&
      p(i)&=\operatorname{softmax}\bigl([u_{t,1},\ldots,u_{t,n};\ a_t]\bigr)_i.
  \end{aligned}
  \label{eq:nlp-extraction-copyre}
\end{equation}
同余均按模3计，$u_{t,i}=\operatorname{selu}(\bm{w}_e^{\top}[\bm{o}_t;\bm{o}^{E}_i])$是复制源位置$i$的分数，
$a_t=\operatorname{selu}(\bm{w}_{\mathrm{NA}}^{\top}\bm{o}_t+b_{\mathrm{NA}})$是NA的分数，selu为缩放指数线性激活，
$\bm{W}_r\in\mathbb{R}^{|R|\times d}$，$\bm{w}_e\in\mathbb{R}^{2d}$。
第一步在$|R|$种关系和NA上选关系；后两步在源句的$n$个位置上选一个位置复制，
依次得到主体和客体，复制客体时用掩码排除刚复制的主体位置。
解码预设生成固定数目的三元组，句中三元组不足时以NA三元组补齐。
这里的复制就是式~\eqref{eq:nlp-analysis-copy}去掉生成门的特例：
关系步相当于$p_{\mathrm{gen}}=1$，而且“词表”换成关系集合；实体步相当于$p_{\mathrm{gen}}=0$，
输出完全由源位置上的分布决定。论文另有每条三元组用一个独立解码器的变体。

在SEO教学句上，参考输出序列为
“主办，复制1，复制6，主办，复制3，复制6”，位置6被复制两次；
EPO句则为“出生地，复制1，复制4，工作地，复制1，复制4”。
训练用教师强制，对每一步取负对数似然；若某条三元组三步的概率依次为0.8、0.7、0.9，
它对损失的贡献为$-\log(0.8\times0.7\times0.9)\approx0.6852$。
生成式输出第一次让同一实体自由参与多条三元组，EPO和SEO都能表示。

代价有三项。其一，每次复制只取一个位置，CopyRE以实体的最后一个词元代表整个实体：
若“海棠大学”被切成“海棠／大学”，模型只复制“大学”，多词实体的边界无法完整恢复，
“海棠大学”与“晨光大学”在输出中也难以区分。其二，三元组本是一个集合，
教师强制却按标注的固定次序计算损失，模型以另一种次序生成同样的集合也会受罚。
其三，三元组数目上限需预先设定。CasRel用首尾指针代替单位置复制，
并把输出改为无次序的条件集合，同时消除了这三项问题。

\subsubsection{CasRel级联主体--关系--客体抽取}
\label{sec:nlp-extraction-casrel}

CasRel先找主体，再在指定主体条件下，按每个关系分别寻找客体\scite{wei2020casrel}
关系不再是两个已选好的实体之间的分类标签，而是从主体到客体集合的映射$f_r(s)$；
映射的值可以为空，也可以有多个客体，因此不需要预设三元组数目，也没有生成次序。
相应地，句子$x$的三元组集合$T$的似然分解为
\[
  p(T\mid x)=\prod_{s\in T}p(s\mid x)
  \prod_{r\in R}p\bigl(\{o:(s,r,o)\in T\}\mid s,x\bigr),
\]
其中$s\in T$指在$T$中作主体的实体，未与$s$构成关系$r$的客体集合为空。
设BERT编码器输出$\bm{h}_i\in\mathbb{R}^{d}$。主体起止头为
$p_i^{\mathrm{sh}}=\sigma(\bm{w}_{\mathrm{sh}}^\top\bm{h}_i+b_{\mathrm{sh}})$与
$p_i^{\mathrm{st}}=\sigma(\bm{w}_{\mathrm{st}}^\top\bm{h}_i+b_{\mathrm{st}})$，
每个位置都是二分类，不要求一句话只有一个主体。对于主体跨度$(a,b)$，取
$\bm{v}_s=(\bm{h}_a+\bm{h}_b)/2\in\mathbb{R}^{d}$，关系$r$的客体头为
\begin{equation}
  p_{i,r}^{\mathrm{oh}\mid s}
  =\sigma\!\left(\bm{w}_{r,\mathrm{oh}}^\top
       (\bm{h}_i+\bm{v}_s)+b_{r,\mathrm{oh}}\right),
  \qquad
  p_{i,r}^{\mathrm{ot}\mid s}
  =\sigma\!\left(\bm{w}_{r,\mathrm{ot}}^\top
       (\bm{h}_i+\bm{v}_s)+b_{r,\mathrm{ot}}\right).
  \label{eq:nlp-extraction-casrel}
\end{equation}
所有$\bm{w}$均为$d$维；关系条件体现在每个关系各自的客体头参数中，共$2|R|$个客体头。
主体标签由全部金标准主体构造，客体标签由“当前主体及当前关系”的全部客体构造。
令$\operatorname{BCE}(y,p)=-y\log p-(1-y)\log(1-p)$，
训练对主体两类边界与选定主体条件下各关系客体两类边界的BCE求和，
原方法每个句子随机抽取一个金标准主体作为条件。

在教学句中，主体起点、终点的正位置均为$\{1,3\}$。给定$s=(1,1)$，
“主办”客体两端的正位置均为6；给定$s=(3,3)$也如此。
假设两次客体起止概率为$(0.9,0.8)$与$(0.85,0.9)$，教学阈值取0.5，
便恢复两条三元组。第一对正边界的损失为$-\log0.9-\log0.8\approx0.3285$；
这只是一对边界的贡献，完整损失还包含负位置及主体预测。
一个主体可以在多个关系头上都有客体，EPO由不同关系头标出同一客体来表示；
多个主体分别作条件，SEO也自然得到表示。

解码先为每个主体起点匹配其后最近的主体终点，再按主体逐关系恢复客体。
这一最近匹配规则简单，却表示不了同一起点对应多个嵌套终点的情形。
更根本的问题是级联：主体漏检会使其全部关系失去后续预测机会；
训练时客体头以金标准主体为条件，推断时却以预测主体为条件，两者分布不同，
这与生成模型的\term{暴露偏差}（exposure bias）同类。推断时每个主体还要单独再算一遍客体头。
下面的单阶段方法正是针对这两点。

\subsubsection{表格填充与TPLinker词元对联合标记}
\label{sec:nlp-extraction-tplinker}

要去掉级联，需要一种输出表示，让全部实体和关系同时预测、彼此不互为条件。
\term{表格填充}（table filling）提供了这种表示\scite{miwa2014table}
对长度为$n$的句子建一张$n\times n$的表，对角格$(i,i)$记录第$i$词的实体标签
（如BILOU加类型），非对角格$(i,j)$记录第$i$词与第$j$词所在实体之间的关系标签，
标签带方向。实体和关系在同一张表上，任何两个实体之间都有自己的格子，SEO不成问题。
原方法按固定次序逐格填写，用特征和束搜索组织决策，后填的格依赖先填的格。
表格表示的限制在于一格只存一个标签：EPO中同一对实体的两种关系落在同一组格子上，无法并存。

TPLinker保留“在词元对上标注”的思路，但给每种关系单独一组标注，
并只标少数几种连接，从而在一步之内解码全部三元组\scite{wang2020tplinker}
它标出三类连接：实体首词元到尾词元（EH--ET），某关系的主体首词元到客体首词元（SH--OH），
以及主体尾词元到客体尾词元（ST--OT）。首首、尾尾两类连接合起来确定了多词实体的两端，
不必像CopyRE那样只取末词。
对$i\leq j$，以
$\bm{g}_{i,j}=\tanh(\bm{W}[\bm{h}_i;\bm{h}_j]+\bm{b})\in\mathbb{R}^{k}$
构造词元对表示，其中$\bm{W}\in\mathbb{R}^{k\times2d}$。
上三角的$n(n+1)/2$个位置按固定顺序铺成一条序列，这种布局称为\term{握手标注}（handshaking tagging）；
它只是共享的索引方式，不是先后执行的推断链。

EH--ET头标记实体首尾，输出“无连接／有连接”。
每个关系另有SH--OH和ST--OT两个头。因为只存上三角，后两类头各含“无连接、正向、反向”三个标签；
反向标签记录主体位置大于客体位置的情形，方向不会因只存上三角而丢失。
共有$2|R|+1$个标记头，各头将$\bm{g}_{i,j}$线性投影到自己的标签空间，
对全部格作交叉熵训练。每种关系各有自己的格子，同一实体对可以在多个关系上都有连接，EPO得以表示。

六词元教学句有21个格。实体头的正格为$(1,1)$、$(3,3)$、$(6,6)$；
“主办”的首部连接和尾部连接均在$(1,6)$、$(3,6)$取正向标签。
推理先由实体头恢复跨度集合$E$。对每个关系$r$，枚举
$s=(a,b),o=(c,d)\in E$，仅当SH--OH给出有向$a\to c$，
且ST--OT给出同方向的$b\to d$时，才恢复$(s,r,o)$。
本例得到两个共享客体的三元组。若只有$(1,6)$的首部连接而没有对应的尾部连接，
就不输出这条关系。

首部和尾部两类连接为什么都需要，可以看一个多词实体。设分词为“海棠／大学／主办／读书会”，
实体头给出$(1,2)$“海棠大学”，同时给出嵌套的$(1,1)$“海棠”（例如标注规范要求标出地名）。
若只有首部连接$1\to4$，无法判断主办者是$(1,1)$还是$(1,2)$；
尾部连接$2\to4$把主体的末端定在位置2，才唯一确定主体为“海棠大学”。
首尾两端同时约束，多词实体和嵌套实体都能完整恢复，这正是CopyRE只复制末词时缺少的信息。

与CasRel相比，TPLinker的所有标签都只以输入为条件，训练与推断完全一致，一次前向即得全部结果。
代价是输出规模：每个关系两个$n(n+1)/2$格的头，总计$O(|R|n^2)$个分类，
长句和大关系集合下显存与计算都显著增加。三类连接独立预测，
在实体边界重复、重叠密集的句子中仍可能错误组合；
实体边界、关系连接和完整三元组的正确率应分别统计。

四种方法每一步扩大了可表示的结构，也增加了解码步骤或输出规模。联合建模并不自动带来更高的准确率：
第~\ref{sec:nlp-extraction-pure}小节的PURE表明，实体信息在输入层注入时，分开训练的流水线同样有竞争力；
选择还取决于实体与关系标注是否齐全，以及句长和关系数目带来的计算代价。

\subsection{文档级关系抽取}
\label{sec:nlp-extraction-docre}

句内方法默认关系的证据在一句话里。DocRED的统计显示，
维基百科文档中相当一部分关系需要跨句才能判断\scite{yao2019docred}
\term{文档级关系抽取}（document-level relation extraction）的输入是整篇文档和其中的实体，
每个实体有若干提及，任务是对每个有序实体对输出关系集合，并可同时给出支持该关系的证据句。
构造三个教学句：$S_1$“海棠大学设有机器学习中心”；$S_2$“该中心负责读书会的筹备”；
$S_3$“海棠楼201配有投影设备”。海棠大学与读书会从未同句出现，若本体与标注规范允许
“所属机构经下属单位参与筹备”这类推理，就要连接$S_1$与$S_2$，并经过“该中心”与“机器学习中心”的共指；
若规范只标直接陈述的关系，就不应补出这条关系。

由此可以列出文档级方法要解决的问题：怎样把一个实体的多个提及聚合为实体表示；
怎样建模经由第三个实体的跨句推理；一个实体对可能有多个关系，也可能没有，阈值怎样决定；
证据句怎样利用。GAIN显式构造文档图，在图上推理；ATLOP不建图，
用预训练编码器的注意力找实体对的上下文，并以可学习的阈值处理多标签；
EIDER在ATLOP的基础上加入证据句预测，并融合原文与证据两路推断。
三种方法都以文档中的实体与提及已经给定为前提。

\subsubsection{文档图推理：GAIN}
\label{sec:nlp-extraction-gain}

跨句推理可以看成图上的路径：海棠大学与机器学习中心在$S_1$中相连，
机器学习中心经共指到$S_2$又与读书会相连。若把这些连接显式建成图，
图神经网络的消息传递就能把信息沿路径送到原本相距很远的实体之间。
GAIN为此构造两层图\scite{zeng2020gain}

第一层是\emph{异构提及图}（heterogeneous mention-level graph），
节点为全部实体提及，外加一个文档节点。三类边各有含义：
同一实体的不同提及之间连\emph{同实体边}，体现共指；
不同实体但出现在同一句中的提及之间连\emph{同句边}，体现局部交互；
每个提及都与文档节点相连，使任意两个提及至多两步可达。
编码器输出的词元表示拼接实体类型嵌入和实体编号嵌入后，取提及内词元的平均作为节点初值$\bm{h}^{(0)}_u$，文档节点的初值取编码器给出的文档表示。
图卷积按边的类型使用不同参数：
\begin{equation}
  \bm{h}^{(l+1)}_u=\sigma\Bigl(\sum_{k\in\mathcal{K}}\sum_{v\in N_k(u)}
    \bm{W}^{(l)}_k\bm{h}^{(l)}_v+\bm{b}^{(l)}_k\Bigr),
  \label{eq:nlp-extraction-gain-gcn}
\end{equation}
$\mathcal{K}$是三种边类型，$N_k(u)$是经类型$k$的边与$u$相连的邻居。
这种按关系类型区分参数的图卷积可参考\BookChapterRef{chap:graph-networks}{卷五的图神经网络章}；
此处只需知道，每层让一个节点汇聚各类邻居的信息，$L$层后信息可传播$L$步。
最终的提及表示拼接各层输出，$\bm{m}_u=[\bm{h}^{(0)}_u;\ldots;\bm{h}^{(L)}_u]$。

第二层是\emph{实体级图}（entity-level graph）。同一实体的提及表示取平均得到实体表示$\bm{e}_i$；
两个实体只要有提及出现在同一句中就连边，边表示
$\bm{e}_{i,j}=\sigma(\bm{W}_q[\bm{e}_i;\bm{e}_j]+\bm{b}_q)$。
判断实体对$(e_h,e_t)$时，GAIN考察经过第三个实体$e_o$的两跳路径，路径表示为
$\bm{p}^{o}_{h,t}=[\bm{e}_{h,o};\bm{e}_{o,t};\bm{e}_{t,o};\bm{e}_{o,h}]$，
再以实体对为查询对路径做注意力：
\[
  \beta_o=\operatorname{softmax}_o\bigl(\sigma([\bm{e}_h;\bm{e}_t]^{\top}\bm{W}_l\,\bm{p}^{o}_{h,t})\bigr),
  \qquad
  \bm{p}_{h,t}=\sum_{o}\beta_o\,\bm{p}^{o}_{h,t}.
\]
分类特征拼接两端实体、它们的差与积、文档节点表示和路径表示，
$\bm{I}_{h,t}=[\bm{e}_h;\bm{e}_t;|\bm{e}_h-\bm{e}_t|;\bm{e}_h\odot\bm{e}_t;\bm{m}_{\mathrm{doc}};\bm{p}_{h,t}]$，
经两层前馈网络对每个关系输出sigmoid概率，用二元交叉熵训练，按开发集选定的阈值输出。
计算量上，提及图的节点数为提及数加一，图卷积的代价与边数成正比；
路径注意力对每个实体对枚举与两端都相连的中间实体，最多$M-2$个，
全部$M(M-1)$个实体对合计为$O(M^3)$，实体多的长文档中这一项不可忽略。

图~\ref{fig:nlp-extraction-gain}展示教学文档上的两层图。提及图中，
$m_2$与$m_3$由同实体边相连，$m_1$与$m_2$、$m_3$与$m_4$分别由同句边相连。
实体图中，海棠大学$e_1$与读书会$e_3$之间没有边，但存在经$e_2$的路径；
路径注意力把这条路径的表示交给$(e_1,e_3)$的分类器。$S_3$中的海棠楼201与三者都不同句，
只经文档节点相连，在路径注意力中不构成两跳路径。

\begin{figure*}[htbp]
  \centering
  % GAIN 的两层图：左为异构提及图（同实体边、同句边、文档边），右为实体级图与两跳路径。
\begin{tikzpicture}[
  x=1mm,y=1mm,figure picture,
  mention/.style={figure node,minimum height=9mm,minimum width=13mm,
    draw=FigureInput,fill=FigureInputFill,font=\figurefont\scriptsize},
  docnode/.style={figure node,minimum height=7mm,minimum width=14mm,
    draw=FigureModel,font=\figurefont\scriptsize},
  entity/.style={figure node,minimum height=9mm,minimum width=15mm,
    draw=FigureInput,fill=FigureInputFill,font=\figurefont\scriptsize},
  intra/.style={draw=FigureStroke,line width=1.1pt},
  inter/.style={draw=FigureStroke,line width=.8pt,dash pattern=on 3pt off 1.5pt},
  docedge/.style={draw=FigureMuted,line width=.6pt,dash pattern=on .8pt off 1.2pt},
  path/.style={-{Stealth[length=1.8mm,width=1.2mm]},draw=FigureOutput,line width=1.1pt},
  query/.style={draw=FigureMuted,line width=.8pt,dash pattern=on 2pt off 2pt},
  note/.style={font=\figurefont\scriptsize,text=FigureMuted,align=center}
]
  \figurePanel{0}{0}{80}{48}{异构提及图}{FigurePanel}
  \figurePanel{88}{0}{60}{48}{实体级图与路径推理}{FigurePanel}

  % 提及图
  \node[docnode] (d) at (40,32) {文档节点};
  \node[mention] (m1) at (10,15) {$m_1$\\海棠大学};
  \node[mention] (m2) at (30,15) {$m_2$\\机器学习\\中心};
  \node[mention] (m3) at (50,15) {$m_3$\\该中心};
  \node[mention] (m4) at (70,15) {$m_4$\\读书会};
  \foreach \m in {m1,m2,m3,m4}{\draw[docedge] (d) -- (\m);}
  \draw[inter] (m1) -- (m2);
  \draw[intra] (m2) -- (m3);
  \draw[inter] (m3) -- (m4);
  \node[note] at (20,4.5) {句$S_1$};
  \node[note] at (60,4.5) {句$S_2$};

  % 实体图
  \node[entity] (e1) at (100,14) {$e_1$\\海棠大学};
  \node[entity] (e2) at (118,30) {$e_2$\\机器学习中心};
  \node[entity] (e3) at (136,14) {$e_3$\\读书会};
  \draw[path] (e1) -- (e2);
  \draw[path] (e2) -- (e3);
  \draw[query] (e1) -- node[figure edge label,font=\figurefont\scriptsize] {?} (e3);
  \node[note] at (118,4.5) {目标实体对$(e_1,e_3)$经$e_2$的两跳路径};
\end{tikzpicture}

  \caption{GAIN在教学文档“$S_1$海棠大学设有机器学习中心；$S_2$该中心负责读书会的筹备”上的两层图。
  左：异构提及图，粗实线为同实体边（“机器学习中心”与“该中心”共指），
  虚线为同句边，点线为各提及到文档节点的边。右：实体级图，同实体提及合并为实体节点，
  同句共现的实体之间连边；目标实体对$(e_1,e_3)$没有直接边，其关系由经$e_2$的两跳路径表示支持。}
  \label{fig:nlp-extraction-gain}
\end{figure*}

显式图的好处是推理路径可见，跨句信息有明确的传递通道。代价在于边由启发式规则决定：
共指是否正确、同句是否意味着相关，都直接影响图的结构，规则漏掉的连接无法由模型补上。
LSR对此的回应是不预先规定边，而把文档图当作潜在结构，用结构化注意力推断边的权重，
并在多轮中迭代细化\scite{nan2020lsr}
它减少了对启发式规则的依赖，却使图的含义更难解释，计算也更重。

\subsubsection{ATLOP文档实体对预测}
\label{sec:nlp-extraction-atlop}

ATLOP的出发点是：预训练Transformer的自注意力已经在全部词元之间建立了连接，
可以看成一张带权的全连接图，长距离依赖不一定要靠另建的图来传递
\scite{zhou2021atlop}
它因此不建图，而用三个部件分别处理多提及、实体对上下文和多标签阈值。
设文档编码为$\bm{H}\in\mathbb{R}^{n\times d}$，每个实体提及前后插入标记“*”。
实体$e$有提及标记表示$\bm{h}_{e,1},\ldots,\bm{h}_{e,M_e}$，逐维对数和指数聚合为
\begin{equation}
  [\bm{h}_e]_u=\log\sum_{j=1}^{M_e}\exp[\bm{h}_{e,j}]_u,
  \qquad u=1,\ldots,d.
  \label{eq:nlp-extraction-atlop-entity}
\end{equation}
它是最大值的平滑近似：每一维主要由取值最大的提及决定，其他提及也有少量贡献。
两个同名对象若被误并为一个实体，聚合结果会直接受到污染。

不同实体对需要文档中不同的上下文。设最后一层注意力有$H_{\!a}$个头，
先在每个头内平均同一实体全部提及标记的注意力行，
得到$\bm{A}_s,\bm{A}_o\in\mathbb{R}^{H_{\!a}\times n}$。
\term{局部上下文池化}（localized context pooling）对两个实体都关注的位置赋权：
\begin{equation}
  q_i=\sum_{h=1}^{H_{\!a}}A_{s,h,i}A_{o,h,i},
  \qquad a_i=\frac{q_i}{\sum_{j=1}^{n}q_j},
  \qquad \bm{c}_{s,o}=\bm{H}^{\top}\bm{a}\in\mathbb{R}^{d}.
  \label{eq:nlp-extraction-atlop-context}
\end{equation}
这里先在同一个头内相乘，再对头求和；先把各头平均再相乘会引入不同头之间的交叉项，
不是同一个计算。该权重反映模型对实体对的关注，不是人工标注的证据概率。

融合后令$\bm{z}_s=\tanh(\bm{W}_s\bm{h}_s+\bm{W}_{cs}\bm{c}_{s,o})$，
$\bm{z}_o$类似，四个投影矩阵均为$d\times d$。
将两向量各分为$K$组，每组$d/K$维，关系$r$的分数为
\begin{equation}
  \ell_r=\sum_{g=1}^{K}
  \bm{z}_{s,g}^{\top}\bm{U}_{r,g}\bm{z}_{o,g}+b_r,
  \qquad
  \bm{U}_{r,g}\in\mathbb{R}^{(d/K)\times(d/K)}.
  \label{eq:nlp-extraction-atlop-score}
\end{equation}
分组把每个关系的双线性参数从$d^2$降到$d^2/K$。以$d=768$、$K=12$为例，每组64维，
参数从$768^2=589{,}824$降到$12\times64^2=49{,}152$。主客体用不同的投影，相同的局部上下文不会使关系自动对称。
除实际关系外，模型还输出一个阈值类别TH，其分数$\ell_{\mathrm{TH}}$也依赖当前实体对。

多标签输出需要阈值。对全部实体对用一个固定阈值，忽略了不同实体对的分数尺度不同。
\term{自适应阈值损失}（adaptive thresholding loss）让阈值成为一个可学习的类别。
令$P$为该对的正关系集合，$N=R\setminus P$为负关系集合，目标把正关系推到TH以上、负关系推到TH以下：
\begin{equation}
  \begin{aligned}
    L_{+}&=-\sum_{r\in P}
      \log\frac{\exp\ell_r}
      {\exp\ell_{\mathrm{TH}}+\sum_{r'\in P}\exp\ell_{r'}},\\
    L_{-}&=-\log\frac{\exp\ell_{\mathrm{TH}}}
      {\exp\ell_{\mathrm{TH}}+\sum_{r'\in N}\exp\ell_{r'}},
    \qquad L_{\mathrm{AT}}=L_{+}+L_{-}.
  \end{aligned}
  \label{eq:nlp-extraction-atlop-loss}
\end{equation}
$L_+$的每一项的分母包含全部正类和TH，不是对每个正类与TH分别做二分类。
$P$为空时$L_+=0$，$L_-$仍训练TH。推理取
$\widehat{R}_{s,o}=\{r\in R:\ell_r>\ell_{\mathrm{TH}}\}$，空集表示NA。

\begin{example}[从局部上下文到多关系输出]
\label{ex:nlp-extraction-atlop}
为单独检查计算，使用四个位置、一个注意力头和二维表示。
令$\bm{H}$的四行依次为$(1,0)$、$(0,1)$、$(2,1)$和$(1,2)$，
前两行是主体的两个提及标记表示，后两行是客体的提及标记表示。
主体和客体分别逐维聚合，得到
\[
  \begin{aligned}
    \bm{h}_s&=(\log(1+\mathrm{e}),\log(1+\mathrm{e})),\\
    \bm{h}_o&=(\log(\mathrm{e}+\mathrm{e}^{2}),
                \log(\mathrm{e}+\mathrm{e}^{2})).
  \end{aligned}
\]
两实体平均后的注意力行取
\[
  \bm{A}_s=(0.4,0.3,0.2,0.1),\qquad
  \bm{A}_o=(0.1,0.2,0.3,0.4).
\]
逐位置乘积为$(0.04,0.06,0.06,0.04)$，其元素之和为0.2。
归一化后，权重向量为$\bm{a}=(0.2,0.3,0.3,0.2)$。
再用同一$\bm{H}$得到上下文向量$\bm{c}_{s,o}=(1.0,1.0)$。

用固定参数使融合向量均为$(0.5,0.5)$：取实体投影为零，
上下文投影为$\operatorname{arctanh}(0.5)\bm{I}_2$。
取$K=1$、偏置为零，并令$r_1,r_2,r_3,\mathrm{TH}$的双线性矩阵依次为
$4\bm{I}_2,2\bm{I}_2,-2\bm{I}_2,\bm{0}$，
式~\eqref{eq:nlp-extraction-atlop-score}给出分数$(2,1,-1,0)$。
若$P=\{r_1,r_2\}$，则
$L_+=2\log(\mathrm{e}^{2}+\mathrm{e}^{1}+1)-3\approx1.8152$，
$L_-=\log(1+\mathrm{e}^{-1})\approx0.3133$，总损失约2.1285。
解码同时保留$r_1,r_2$，排除$r_3$。这里的简化参数只用来贯通计算，
不表示训练好的模型会忽略实体表示。
\end{example}

与GAIN相比，ATLOP的取舍很清楚。它不需要共指、同句共现等规则来定义边，
也没有额外的图网络层，实现和训练都更简单；实体对的上下文由预训练注意力直接给出，
不受启发式边的遗漏所限。放弃的是显式的推理结构：经由哪个中间实体、走了哪条路径，
在ATLOP中只隐含在注意力权重里，无法像GAIN的路径注意力那样逐条检查；
需要多步推理且中间实体与两端都相距很远时，单层注意力乘积未必能把它们联系起来。
文档超过编码器长度上限时，还要分段编码再合并注意力与表示。

\subsubsection{EIDER证据增强与推理融合}
\label{sec:nlp-extraction-eider}

ATLOP的局部上下文是模型自己学到的关注，没有直接训练“哪些句子足以支持关系”。
DocRED标注了\term{证据句}（evidence sentence），可以作为额外监督。
EIDER在共享的文档编码器上增加实体对条件的证据预测，并在推理时结合完整文档与证据伪文档的关系分数
\scite{xie2022eider}
设文档有$J$句，第$j$句的表示$\bm{s}_j\in\mathbb{R}^{d}$由句内词元向量逐维对数和指数池化获得。
给定实体对上下文$\bm{c}_{s,o}$，第$j$句是证据的概率为
\begin{equation}
  p_j^{s,o}=\sigma(\bm{s}_j^\top\bm{W}_v\bm{c}_{s,o}+b_v),
  \qquad \bm{W}_v\in\mathbb{R}^{d\times d}.
  \label{eq:nlp-extraction-eider-evidence}
\end{equation}
证据按实体对选取，而不是为每个关系类别分别选一套句子。

设$D^+$为至少有一个正关系的训练实体对集合，$y_j^{s,o}$为该对第$j$句的证据标签，
证据损失为二元交叉熵，与关系损失加权相加：
\begin{equation}
  \begin{aligned}
    L_{\mathrm{evi}}
      &=-\sum_{(s,o)\in D^+}\sum_{j=1}^{J}
      \left[y_j^{s,o}\log p_j^{s,o}
       +(1-y_j^{s,o})\log(1-p_j^{s,o})\right],\\
    L&=L_{\mathrm{AT}}+\lambda L_{\mathrm{evi}}.
  \end{aligned}
  \label{eq:nlp-extraction-eider-loss}
\end{equation}
证据目标只在有正关系的实体对上计算，NA实体对没有可定义的证据。
两个任务共同更新编码器，证据监督让编码器更关注真正支持关系的句子。
没有人工证据标注时，EIDER用启发式规则构造证据：依次尝试两实体同句共现的句子、
含共指提及的句子、以及经一个桥接实体分别与两端相连的句组。规则证据可能包含与目标关系无关的句子，
它是弱监督，不能当作人工核验。

推理分两路。对每个实体对按阈值选出证据句，保持原句序拼成伪文档，
同时重建提及位置以及伪文档位置到原文位置的映射；
关系模型分别在完整文档和伪文档上给出分数$\ell_r^O$与$\ell_r^E$，各自减去各自的TH分数，
再以一个在开发集上拟合的标量$\tau$融合：
\begin{equation}
  S_r^O=\ell_r^O-\ell_{\mathrm{TH}}^O,\qquad
  S_r^E=\ell_r^E-\ell_{\mathrm{TH}}^E,\qquad
  p_r^{\mathrm{fuse}}=\sigma(S_r^O+S_r^E-\tau).
  \label{eq:nlp-extraction-eider-fusion}
\end{equation}
以$p_r^{\mathrm{fuse}}>0.5$判正等价于$S_r^O+S_r^E>\tau$。
完整文档一路保留了所有背景，证据一路去掉了干扰句，两路分数相加让二者互相补充。
这种\term{推理阶段融合}（inference-stage fusion）不需要重新训练分类器。
选出的证据为空或缺少某个实体的提及时，应退回完整文档一路的结果，并记录回退。

沿用上面的教学文档。设三句的证据标签为$(1,1,0)$，
预测概率为$(0.8,0.7,0.2)$，证据损失之和为
$-\log0.8-\log0.7-\log0.8\approx0.8030$。
取0.5为选句阈值，伪文档包含$S_1,S_2$。假设某个允许的筹备关系在两路的“关系、TH”分数
分别为$(0.6,0.2)$和$(0.8,0.1)$，则$S^O=0.4$、$S^E=0.7$。取教学融合阈值$\tau=0.3$，
得到$\sigma(0.8)\approx0.6900$，输出该关系。
若证据预测漏掉$S_1$，伪文档中缺少海棠大学与中心的联系，第二路可能变为$S^E=-0.5$，
融合概率降为$\sigma(-0.4)\approx0.4013$，关系被漏掉。
可见证据一路的收益依赖证据预测的召回；证据遗漏、背景误选与关系误判应分别检查。

表~\ref{tab:nlp-extraction-docre-methods}汇总本小节的文档级方法：从GAIN到ATLOP，跨句信息的传递从显式图转到预训练注意力；
从ATLOP到EIDER，模型关注的上下文从无监督的注意力转到有监督的证据句。

\begin{table*}[htbp]
  \centering
  \small
  \begin{tabularx}{\linewidth}{@{}>{\raggedright\arraybackslash}p{12mm}
    *{3}{>{\raggedright\arraybackslash}X}@{}}
    \toprule
    方法 & 解决的问题 & 核心机制 & 新的代价 \\
    \midrule
    GAIN & 跨句推理缺少传递通道 & 异构提及图上按边类型的图卷积；实体级图上的两跳路径注意力 & 边依赖共指与同句规则；增加图网络层 \\
    LSR & 启发式边可能遗漏或出错 & 以结构化注意力推断潜在文档图，并迭代细化 & 图的含义难解释；计算更重 \\
    ATLOP & 多提及聚合、实体对上下文与多标签阈值 & 对数和指数聚合；同头注意力乘积池化；可学习阈值类TH & 推理路径隐含在注意力中 \\
    EIDER & 关注的上下文未经证据监督 & 实体对条件的证据预测；原文与证据伪文档两路分数融合 & 需证据标注或规则；推断需两次编码 \\
    \bottomrule
  \end{tabularx}
  \caption{文档级关系抽取方法对照。各方法均假设文档中的实体及其提及已经给定；
  “新的代价”指相对于前一行方法新增的依赖或计算。}
  \label{tab:nlp-extraction-docre-methods}
\end{table*}

\subsection{开放关系抽取：ReVerb}
\label{sec:nlp-extraction-reverb}

前面的方法都要求预先定义$R$。\term{开放信息抽取}（open information extraction, OpenIE）
不预设关系类型，直接从句子中抽取“论元—关系短语—论元”的三元组，关系用原文的\term{关系短语}
（relation phrase）表示。它适合类型无法穷举、需要先探索语料再决定本体的场合。
早期开放抽取系统常产生两类坏结果：关系短语不连贯，由不相邻、拼不成意义的词凑成；
或信息不足，例如把“made a deal with”只抽成“made”，丢掉了承载关系含义的名词。
ReVerb用句法约束抑制这两类错误，再用词汇约束排除只对个别论元成立的过度特定短语\scite{fader2011reverb}

句法约束要求关系短语匹配词性模式$V\mid VP\mid VW^{*}P$：
$V$为动词（可带小品词和副词），$P$为介词、小品词或不定式标记，
$W$为名词、形容词、副词、代词或限定词；多个匹配相邻时合并，并取最长匹配。
这一模式覆盖了“hosts”“is held in”“has a meeting with”等形式，又排除了不以动词为核心的片段。
词汇约束要求关系短语在大规模语料中与足够多不同的论元对共现，
论文的阈值为20对，用以排除只对某一对论元成立的过度特定短语。
找到关系短语后，在其左右各取最近的名词短语作论元，并排除关系代词、疑问词和存在句中的“there”等。
最后用逻辑回归给每个抽取打置信度，$p(\text{正确}\mid\bm{\phi})=\sigma(\bm{w}^{\top}\bm{\phi})$，
特征描述句子长度、关系短语形式、论元位置等，训练数据是少量人工判定正误的抽取。

把补充句译写为英语教学句“Haitang University hosts the reading group.”，
关系短语为“hosts”，左侧最近的名词短语为“Haitang University”，右侧为“the reading group”，
得到开放三元组(Haitang University, hosts, the reading group)。
封闭模式可以把“hosts”归一为“主办”；开放结果保留原措辞及其位置，归一需要另外的映射。
ReVerb的规则与词汇统计面向英语，中文需要另行设计词性模式与资源。
开放三元组较短，“does not host”中的否定、时间状语和条件从句若被丢弃，原意就会改变；
这些限定需要下一小节记录中的专门字段承接。

\subsection{关系记录的恢复与核验}
\label{sec:nlp-extraction-relation-records}

模型输出的是分数和词元位置，可用的记录还需要对象身份、规范关系、限定条件和来源。
教学记录可写为$\langle s,r,o,\text{polarity},\text{time},\text{condition},\text{source}\rangle$；
否定关系记为负极性，而不是删除后留下同样的正关系。
算法~\ref{alg:nlp-extraction-relation-records}给出从分数到记录的步骤，
多条记录之间的对齐、合并与更新见第~\ref{sec:nlp-extraction-records}节。
结构检查使用JSON Schema约束字段类型、必填项和枚举值\cite{wright2022jsonschema}，
用SHACL约束图中的类型、属性取值范围和基数\cite{knublauch2017shacl}；
它们检查的是格式与结构，内容是否为原文所支持，要靠证据与身份的逐项核对。

\begin{algorithm}[关系记录的恢复与核验]
  \label{alg:nlp-extraction-relation-records}
  \begin{algorithmic}[1]
    \Require 文档及其版本号；候选三元组及分数、阈值；本体（关系方向、类型约束、基数）；
      别名到规范关系的映射；记录模式
    \Ensure 通过检查的关系记录，以及带拒绝原因的候选
    \For{每个分数超过阈值的候选$(s,r,o)$}
      \State 由词元跨度恢复原文字符区间；重复字符串按主客体与上下文消歧
      \State 开放关系短语按映射归一，同时保存原措辞；被动等表述按映射调整主客体方向
      \State 把$s,o$链接到库中身份或标为未链接提及；检查类型是否满足$r$的定义域与值域
      \State 从证据范围读取否定、时间和条件，填入极性、时间与条件字段；时间按日期规范解析
      \State 记录来源：文档版本、证据句与字符区间
      \State 用记录模式检查字段，用图约束检查类型与基数；不通过则附原因并拒绝
      \State 以证据为前提、记录的文字表述为假设做支持判定，并单独复核身份、数值、否定与条件
    \EndFor
    \State 精确匹配按对象键、规范关系与条件字段比较；语义等价评价另声明允许的别名、时间表达与方向改写
  \end{algorithmic}
\end{algorithm}

算法中的几道检查各管一层：类型兼容能拒绝“日期主办组织”，却不能判断类型兼容的边是否真实；
格式通过也不能抵消内容错误。“未参与主办”与“参与主办”只差几个字却不等价，
而“负责举办”与“主办”是否等价要由预先声明的映射决定。
流水线的候选遗漏、联合模型的连接误组、生成模型的无依据补项，最终都落到同一组对象、
关系、条件与来源的检查上；给定正确实体条件下的关系分数，不能代替从原文开始的端到端结果。
